\documentclass[letterpaper]{article} % DO NOT CHANGE THIS
\usepackage{aaai2027}  % DO NOT CHANGE THIS

\nocopyright

% AAAI-compatible packages.
\usepackage[hyphens]{url}
\usepackage{graphicx}
\usepackage{natbib}
\usepackage{caption}
\usepackage{booktabs}
\usepackage{amsmath}
\usepackage{amssymb}
\usepackage{algorithm}
\usepackage{algorithmic}

\urlstyle{rm}
\def\UrlFont{\rm}

\title{Solver-Integrated Adversarial Attacking and Training of Neural Operators}

\author{
  Yifei Sun\textsuperscript{\dag}
}
\affiliations{}

\begin{document}

\maketitle
\begingroup
\renewcommand{\thefootnote}{\fnsymbol{footnote}}
\footnotetext[2]{\textsuperscript{\dag} Corresponding author. Email: \texttt{yifeisun@umich.edu}.}
\endgroup

\begin{abstract}
Neural operators are commonly evaluated as fast surrogates for numerical
solvers, yet their generalization and robustness are often measured without
fully using the solver that defines the underlying PDE solution map. This paper
studies the solver-integrated setting in which a neural operator and a
classical solver act as two operators on the same function space. We distinguish
static generalization, measured by the discrepancy between the model and solver
at a fixed input, from local robustness, measured by how this discrepancy
changes under norm-bounded perturbations of the input function. From this view,
we define an error operator and compare three robustness notions: adversarial
maximization of the final solver-model error, local operator sensitivity through
a Jacobian-like linearization, and a first-order Jacobian-error proxy that
predicts the loss-increasing perturbation direction. The resulting analysis
shows why attacks that only enlarge the neural operator output, or that keep
the solver target fixed, can be inefficient for regression PDE problems: the
model and solver may move together, leaving the true error nearly unchanged. We
therefore propose solver-integrated adversarial attacks and training objectives
that optimize the true perturbed discrepancy. Experiments on Burgers' equation,
Darcy flow, and Navier--Stokes settings show that stronger solver integration
finds more informative failure modes and improves both generalization and
robustness, while exposing the computational cost of differentiating through
time-dependent solvers.
\end{abstract}

\section{Introduction}

Neural operators, including Fourier Neural Operators and DeepONets, have
become a prominent approach for learning PDE solution maps from data
\citep{li2021fourier,lu2021deeponet}. Their reliability under perturbations is
now an active question. Classical adversarial learning studies motivate
worst-case perturbation training and evaluation
\citep{goodfellow2015explaining,madry2018towards,zhang2019trades,croce2020autoattack}.
Adesoji and Chen evaluated adversarial robustness for
Fourier Neural Operators by comparing perturbed predictions with PDE solver
outputs \citep{adesoji2022evaluating}. StablePDENet formulates operator
learning as a stability-aware min-max problem and uses adversarial training to
improve fidelity under worst-case input perturbations
\citep{huang2026stablepdenet}. Related work also includes adversarial-gradient
sample movement for scientific machine learning and operator learning
\citep{ouyang2025rams}, white-box residual-region refinement for PINNs
\citep{shi2025wbar}, sparse gradient-free attacks on neural-operator
digital twins \citep{roy2026vulnerabilities}, and active-learning plus
denoising defenses for robust neural operators \citep{roy2026beyond}.
Uncertainty-driven active operator learning is another adjacent line:
Winovich et al. use predictive uncertainty to select PDE instances for
additional solver labeling, while MRA-FNO chooses both input functions and
simulation resolutions through a utility/cost acquisition rule
\citep{winovich2025active,li2024mrafno}. These methods reduce the cost of
solver-generated datasets and improve clean generalization, but their
acquisition objectives are not worst-case perturbation attacks and do not
maximize a solver-model discrepancy.
Another adversarial line is generative rather than worst-case-perturbation
based: Generative Adversarial Neural Operators learn function distributions
\citep{rahman2022gano}, and adversarially trained neural operators have been
used to recover sharper turbulent-flow structures and energy spectra
\citep{oommen2025turbulent}.

This paper is also connected to the DeepONet and physics-informed neural
operator line of work associated with Karniadakis and collaborators
\citep{raissi2019physics,li2021pino,lu2021faircomparison,goswami2022physicsinformed}.
That literature highlights operator learning, PDE residual information, and
noise sensitivity, but the role of a numerical solver as a differentiable
participant in both the forward and backward attack loop remains
underexplored. Our focus is the solver-integrated regression setting: the
adversary should enlarge the true perturbed discrepancy between the learned
operator and the solver, rather than only moving the learned model output or
using a fixed target.

\section{Problem Setting}

Let $\mathcal{X}$ and $\mathcal{Y}$ be Banach or Hilbert spaces of functions,
for example spaces of initial conditions, coefficient fields, source terms,
or solution fields on a spatial or space-time domain. An input is a function
$a\in\mathcal{X}$ and an output is a function in $\mathcal{Y}$. We consider a
learned neural operator
\begin{equation}
  F_\theta:\mathcal{X}\to\mathcal{Y}
\end{equation}
and a numerical solver or oracle
\begin{equation}
  G:\mathcal{X}\to\mathcal{Y}
\end{equation}
for the same PDE solution map. All perturbation budgets, energies, and
errors are therefore norms of functions unless explicitly stated otherwise.
For a norm $\|\cdot\|_{\mathcal{X},p}$ on the input space, the local
perturbation set around $a$ is the function-space ball
\begin{equation}
  \mathcal{B}_{p,\varepsilon}(a)
  =
  \{a+\eta:\eta\in\mathcal{X},\ \|\eta\|_{\mathcal{X},p}\le\varepsilon\}.
\end{equation}

The finite-dimensional viewpoint is only the discretized realization of this
operator-level formulation. After choosing a mesh or basis, $a$ and
$F_\theta(a)$ are represented by vectors, and the operators become nonlinear
maps between high-dimensional Euclidean spaces. In that representation,
function norms are approximated by weighted vector norms, Frechet derivatives
become Jacobian matrices, adjoint derivatives become transposed Jacobians, and
induced operator norms become matrix norms such as spectral norms when
$p=q=2$.

\section{Solver-Integrated Metrics}

The model-solver error operator is the map
\begin{equation}
  E_\theta := F_\theta - G,\qquad E_\theta:\mathcal{X}\to\mathcal{Y}.
\end{equation}
For a particular input function $u\in\mathcal{X}$, the corresponding error
field is the value of this operator,
\begin{equation}
  e_\theta[u] := E_\theta(u)=F_\theta(u)-G(u)\in\mathcal{Y}.
\end{equation}
Thus $E_\theta$ is the operator, $e_\theta[u]$ is an output function, and the
vector of sampled error values appears only after discretization. Static
generalization at an input function $u$ is
\begin{equation}
  \mathcal{L}_{\mathrm{gen}}(u)
  =
  \|e_\theta[u]\|_{\mathcal{Y},q}
  =
  \|F_\theta(u)-G(u)\|_{\mathcal{Y},q}.
\end{equation}
It measures the model-solver discrepancy at a fixed function. Local
robustness instead asks how large the same discrepancy can become in a
neighborhood of that function:
\begin{equation}
  R_{\varepsilon,p,q}(u)
  =
  \sup_{\|\eta\|_{\mathcal{X},p}\le\varepsilon}
  \|e_\theta[u+\eta]\|_{\mathcal{Y},q}.
\end{equation}

If $E_\theta$ is Frechet differentiable at $u$, then for small perturbations
\begin{equation}
  e_\theta[u+\eta]
  =
  e_\theta[u]+D E_\theta[u](\eta)+o(\|\eta\|_{\mathcal{X}}),
\end{equation}
where $D E_\theta[u]:\mathcal{X}\to\mathcal{Y}$ is a bounded linear operator.
This gives the local operator-sensitivity metric
\begin{equation}
  \|D E_\theta[u]\|_{p\to q}
  =
  \sup_{\|\eta\|_{\mathcal{X},p}\le 1}
  \|D E_\theta[u](\eta)\|_{\mathcal{Y},q}.
\end{equation}
In a discretized Hilbert-space setting with $p=q=2$, this is the spectral
norm of the error Jacobian. It measures the largest local change of the
error, not necessarily the largest final error.

For the squared Hilbert norm
$\Phi(u)=\frac{1}{2}\|e_\theta[u]\|_{\mathcal{Y}}^2$, the first variation is
\begin{align}
  D\Phi[u](\eta)
  &=
  \langle e_\theta[u],D E_\theta[u](\eta)\rangle_{\mathcal{Y}}
  \nonumber\\
  &=
  \langle D E_\theta[u]^*e_\theta[u],\eta\rangle_{\mathcal{X}},
\end{align}
where $D E_\theta[u]^*:\mathcal{Y}\to\mathcal{X}$ is the adjoint derivative.
Thus the operator-level analogue of the finite-dimensional
Jacobian-error vector is
\begin{equation}
  D E_\theta[u]^*e_\theta[u],
\end{equation}
which discretizes to $J_E(x)^\top e_\theta(x)$. This quantity incorporates both the
local linearization and the current error direction, so it is a first-order
proxy for the perturbation that increases the final solver-model discrepancy.

This notation also separates three attack targets:
\begin{align}
  L_1(\eta) &= \|F_\theta(u+\eta)-F_\theta(u)\|_{\mathcal{Y},q},\\
  L_2(\eta) &= \|F_\theta(u+\eta)-G(u)\|_{\mathcal{Y},q},\\
  L_3(\eta) &= \|F_\theta(u+\eta)-G(u+\eta)\|_{\mathcal{Y},q}
             =\|e_\theta[u+\eta]\|_{\mathcal{Y},q}.
\end{align}
Classical fixed-label attacks are closest to $L_2$-type objectives because
the reference target is held fixed. In PDE regression, however, the solver
output is itself a function of the perturbed input, so the solver-consistent
objective is $L_3$.

\section{Solver-Integrated Adversarial Attack}

Write the attack objective and optimization procedure here.

\section{Solver-Integrated Adversarial Training}

Write the training objective and data-generation procedure here.

\section{Experiments}

Write the Burgers, Darcy flow, and Navier--Stokes experiments here.

\section{Conclusion}

Write the conclusion here.

\bibliography{references}

\end{document}
